\chapter{Gestión y metodología de trabajo}\label{cap:gestion}

El autor organizó SavePoint en siete fases, desde catálogo y biblioteca hasta la comparación experimental y la preparación de demostración. La planificación se conserva en \texttt{PROJECT.md}, \texttt{REQUIREMENTS.md}, \texttt{ROADMAP.md}, \texttt{STATE.md}, planes, resúmenes y verificaciones. Esta cadena permite enlazar una decisión con un cambio, una prueba y su evidencia; la matriz de trazabilidad conserva los detalles sin sobrecargar la memoria principal.

El autor aplicó una metodología de ingeniería asistida por agentes, gobernada mediante GSD y organizada conceptualmente con Obsidian. El ciclo de discutir, investigar, planificar, ejecutar y verificar conserva contexto entre tareas. Los agentes se emplearon como herramientas para planificar, investigar, revisar, auditar y generar propuestas; el autor revisó, modificó o rechazó esas propuestas y asumió todas las decisiones y la responsabilidad académica. El vault enlaza conceptos, fases y ADR con fuentes canónicas, pero no sustituye tests, hashes ni artefactos versionados.

La metodología reduce pérdida de contexto y facilita reanudación, pero no garantiza corrección. Depende de especificaciones claras, requiere revisión humana y puede dejar documentación desactualizada. \todo[inline]{PENDIENTE: incorporar horas, calendario y presupuesto con el registro que el autor considere defendible.}
